\documentclass[10pt]{article}

% graphicx package, useful for including eps, pdf, and raster graphics
\usepackage{graphicx}
\DeclareGraphicsExtensions{.png,.png,.jpg,.pdf}

% basic packages
\usepackage{fancyhdr}
\usepackage{color}
\usepackage{parskip}
\usepackage{float}
\usepackage{microtype}
\usepackage{url}
\usepackage{hyperref}
\usepackage{booktabs}
\usepackage{makecell}
\usepackage{multirow}
\usepackage{pbox}
\usepackage{enumitem}
\usepackage{xeCJK}

% fonts
\usepackage{fontspec}
\setmainfont{DroidSans}[
  Extension   = .ttf,
  UprightFont = *,
  BoldFont    = *-Bold ]
\setmonofont{DroidSansMono}[
  Extension   = .ttf,
  UprightFont = * ]
\setCJKmainfont{NotoSansCJK}[
  Extension   = .ttc,
  UprightFont = *-Regular,
  BoldFont    = *-Bold,
  Language    = Chinese Traditional ]

\hypersetup{colorlinks,urlcolor=blue}
\urlstyle{same}

% page layout
\usepackage[top=1in, bottom=1in, left=0.75in, right=0.75in, includefoot, heightrounded]{geometry}

% Useful for seeing the layout as you tweak it.
%\usepackage{showframe}

\setlength{\headheight}{1.1in}
\pagestyle{fancy}
\fancyhf{}

\lhead{\pbox[b][1in][c]{\textwidth}{
    \small
    \textbf{检测执行方：}\\
    Northwest Genomics Center\\
    University of Washington School of Medicine\\
    Genome Sciences Box 355065\\
    3720 15th Ave NE\\
    Seattle, WA 98195-5065}}
\rhead{\includegraphics[height=0.75in,trim=0 -25 0 0]{logo}}

\newcommand{\PageLine}{\rule{\textwidth}{0.25mm}}

\begin{document}
\sloppy % prevent overfilled lines which occur with long URLs

\begin{description}[font=\normalfont,align=left,labelwidth=6em]
\item [\textbf{\VAR{pat_name|e}}]
\item [\textbf{出生日期：}] \VAR{birth_date|e}
\item [\textbf{標本：}]
  %- if swab_type == "ans"
  鼻腔拭子,
  %- elif swab_type == "mtb"
  鼻腔中部，
  %- else
  \errmessage{Invalid swab_type "\VAR{swab_type}"
  %- endif

  自我採集

\item [\textbf{採集條碼：}] \VAR{qrcode|e}
\item [\textbf{採集日期：}] \VAR{collect_ts|e}
\item [\textbf{報告日期：}]
  %- if status_code == "never-tested"
  沒有進行檢測
  %- else
  \VAR{result_ts|e}
  %- endif
\end{description}

\PageLine

\begin{center}
\Large
\textbf{研究結果 – 鼓勵您與您的醫療保健提供者分享此結果}
\end{center}

\bigskip

\large \underline{\textbf{關於本檢測}}

\begin{itemize}
\item

  COVID-19 是一種名為 SARS-CoV-2 的\textbf{病毒}引發的\textbf{疾病}。
  本研究檢測會在從您的鼻子裡採集的樣本中尋找這種病毒，以確定您是否患有 COVID-19。
  這種檢測的技術術語為「逆轉錄聚合酶鏈反應」，簡稱「RT-PCR」。

\end{itemize}

\bigskip

\large \underline{\textbf{重要提示}}

查看您的研究檢測結果前，您應該確認以下內容：

\begin{itemize}
\item

  您收集的\textbf{樣本類型}和\textbf{實驗室檢測}是\textbf{研究}的一部分內容，
  本報告介紹了您的\textbf{研究檢測結果。}

\item
  %- if staff_observed == "f"
  您自己（或為他人）採集樣本並貼上標籤後，透過快遞將其送回實驗室。
  它不是由護士或醫生等有執照的醫療保健專業人員採集、貼標籤或運送的

  %- elif staff_observed == "t"
  您自行（或為他人）採集樣本，由快遞員或其他類型的非醫療保健專業人員貼上標籤並將其送回實驗室。
  它不是由護士或醫生等有執照的醫療保健專業人員採集或運送的。

  %- else
  \errmessage{Unknown value: staff_observed = "\VAR{staff_observed}"}

  %- endif

  這可能會影響您的檢測結果，醫療保健提供者可能希望再次對您進行檢測。

\item

  參與本研究並不能代替常規醫療護理。如果您感覺很不舒服或擔心您的檢測結果，
  您應該聯繫醫生或醫療保健提供者。

\end{itemize}

\bigskip
\PageLine

\large \underline{\textbf{您的 SARS-CoV-2 (COVID-19) 研究檢測結果}}

%- if status_code == "negative"
未檢出（陰性）\\
%- elif status_code == "positive"
檢出（陽性）\\
%- elif status_code == "inconclusive"
不確定\\
%- elif status_code == "never-tested"
無法進行檢測\\
%- else
\VAR{0/0} %# poor man's raise
%- endif

\PageLine
\bigskip

\large \underline{\textbf{瞭解您的結果}}

此表給出了您的檢測可能顯示的所有結果，並描述了每個結果的含義。
您可以列印或透過電子郵件將此\textbf{研究檢測結果}與您的醫生或醫療保健提供者共享

不管您的檢測結果如何，如果您感覺不舒服，您應該待在家裡，避免和其他人在一起。
\textbf{如果您一直有症狀，您病得很重或擔心您自己的疾病，或對您的研究檢測結果有疑問，
您應該聯繫醫生或醫療保健提供者。}

\begin{itemize}


\item

  \textbf{若您的檢測結果顯示「未檢出（陰性）」}，即表示我們實驗室
  未在您寄送給我們的樣本中發現引起 COVID-19 的病毒。沒有什麼檢測是完美的，
  即使您的檢測結果為陰性，您也有些許可能性感染病毒。

\item

  \textbf{若您的檢測結果顯示「檢出（陽性）」}，即表示我們實驗室
  在您寄送給我們的樣本中發現引起 COVID-19 的病毒。您應該待在家中，避免接觸其他人，
  尤其是您感到不舒服的時候。

  您的檢測結果將與華盛頓州和當地衛生部門共享。您不需要自己報告結果。衛生部門可能會打電話給您，
  詢問您的病情以及您在生病時所接觸人員的情況。

\item

  \textbf{若您的檢測結果顯示「不確定」}，即表明我們實驗室對您的檢測結果沒有把握。
  沒有什麼檢測是完美的，當我們在您的樣本中發現少量的 SARS-CoV-2，但數量不足以讓我們說您患有
  COVID-19 時，就會出現不確定的結果。

\item

  \textbf{若您的檢測結果顯示「無法進行檢測」}，即表示 您的樣本中沒有足夠的物質供我們檢測，
  或者我們進行檢測時出現錯誤。由於我們無法檢測您的樣本，所以無法發現您是否攜 帶導致 COVID-19
  的病毒。

\end{itemize}

\bigskip

\large \underline{\textbf{資源}}

如果您想瞭解更多關於 COVID-19 的資訊，請查看公共衛生合作夥伴提供的這些資源:

\begin{itemize}
\item
  Public Health \textemdash Seattle \& King County（西雅圖和金縣公共衛生）\\
  \url{https://www.kingcounty.gov/depts/health/covid-19/languages/chinese.aspx}

\item

  Washington State Department of Health（華盛頓州衛生部）\\
  \url{https://www.doh.wa.gov/Emergencies/NovelCoronavirusOutbreak2020COVID19/ChineseSimplified}

\item
  美國 Centers for Disease Control and Prevention（CDC, 疾病控制和預防中心）\\
  \url{https://chinese.cdc.gov/coronavirus/2019-ncov/index.html?CDC\_AA\_refVal=https%3A%2F%2Fwww.cdc.gov%2Fcoronavirus%2F2019-ncov%2Findex-Chinese.html}

\end{itemize}

If you would like to speak to someone on the phone, please call one of the
following:

\begin{itemize}

\item

  \textbf{與 COVID-19 有關的醫療問題：} 早上 8 點至晚上 7 點，可撥打
  \textbf{206-477-3977} 聯繫 King County Call Center（金縣呼叫中心）

\item

  \textbf{關於 COVID-19 的非醫學問題，包括合規性和企業相關問題:} 週一至週五早上 8:30
  至下午 4:30，可撥打 \textbf{206-296-1608} 聯繫 King County COVID-19 Business and Community
  Information Line（金縣 COVID-19 企業和社區資訊熱線）

\item

  \textbf{關於華盛頓州 COVID-19 的一般問題:} 早上 6 點至晚上 10 點，可撥打
  \textbf{1-800-525-0127} 聯繫 Washington State Novel Coronavirus Call Center
  （華盛頓州新型冠狀病毒呼叫中心）

\end{itemize}

\bigskip

\large \underline{\textbf{方法和局限性}}

\textbf{方法}\\
本實驗使用逆轉錄 PCR 檢測 SARs-CoV-2。該實驗使用病毒片段 ORf1B 和 S 的擴增，並檢測人類
RnaseP 擴增對照。該檢測由 University of Washington 的 Northwest Genomics Center
開發並確定其性能特徵。該實驗室獲得 WA Department of Health醫學檢驗站的許可，
具備進行高複雜性臨床實驗室檢測的資格。

\textbf{局限性}

\begin{itemize}

\item

  這些結果只是研究的一部分內容。

\item
  %- if staff_observed == "f"
  您自己（或為他人）採集樣本並貼上標籤後，透過快遞將其送回實驗室。
  它不是由護士或醫生等醫療保健專業人員採集、貼標籤或運送的。

  %- elif staff_observed == "t"
  您自行（或為他人）採集樣本，由快遞員或其他類型的非醫療保健專業人員貼上標籤並將其送回實驗室。
  它不是由護士或醫生等有執照的醫療保健專業人員採集或運送的。

  %- else
  \errmessage{Unknown value: staff_observed = "\VAR{staff_observed}"}

  %- endif
  這可能會影響您的檢測結果，醫療保健提供者可能希望再次對您進行檢測。

\item

  由於此報告基於一項研究資料，因此您的結果可能無法診斷由 SARS-CoV-2 感染引起的 COVID-19。

\end{itemize}

\end{document}
